We propose a deployable, data-driven fan vote inference model. The key design principle is to \textbf{treat audience votes as a latent signal and infer vote shares from observed judges' scores and elimination outcomes under explicit rule constraints.}

\subsection{Overview and Modeling Philosophy}\label{subsec:model_philosophy}

\textbf{Fan votes are a latent signal and must be inferred indirectly.}
In DWTS, eliminations depend on a combination of judges' scores and audience votes, yet only judges' scores and realized eliminations are observed.
We therefore model weekly audience support as a latent vote-share vector over the active set.
Because many vote configurations can be consistent with the same elimination outcome, the inference problem is generally underdetermined; our approach resolves this by
(i) enforcing elimination-consistency constraints implied by the show's rule mechanics and
(ii) selecting a stable solution via temporal smoothness regularization.\cite{benning2018regularization,kaipio2005statistical}
Special episodes (e.g., no-elimination or irregular-format weeks) are handled explicitly through slack variables rather than forcing exact compliance, and these slacks become a diagnostic for identifiability and model mismatch.

\subsection{Main Inference Model}\label{subsec:inference_model}
\textbf{Decision variables.}
For each season $s$, week $t$, and active contestant $i\in A_{s,t}$, let $p_{i,t}\ge 0$ denote the unobserved fan vote share, with $\sum_{i\in A_{s,t}}p_{i,t}=1$.
Let $J_{i,t}$ be the observed total judges' score (sum across judges with missing values skipped) and define the judges' share
$q^J_{i,t}=J_{i,t}\big/\sum_{r\in A_{s,t}}J_{r,t}$.

\textbf{Percent-scheme combined score (primary inference target).}
Under the percent-based aggregation, we define the combined score
\begin{equation}
S^{(P)}_{i,t}=\alpha q^{J}_{i,t}+(1-\alpha)p_{i,t},
\end{equation}
where $\alpha\in[0,1]$ controls the relative weight of judges versus fans (default $\alpha=0.5$; later stress-tested).

\textbf{Elimination-consistency constraints with slack.}
Let $E_{s,t}\subseteq A_{s,t}$ denote the observed eliminated set at the end of week $t$ with $m_{s,t}=|E_{s,t}|$.
We introduce a nonnegative slack $\xi_t\ge 0$ and impose
\begin{equation}
S^{(P)}_{e,t}\le S^{(P)}_{i,t}+\xi_t,\quad \forall e\in E_{s,t},\ \forall i\in A_{s,t}\setminus E_{s,t}.
\end{equation}

When the assumed rule and data are perfectly compatible, the minimum feasible $\xi_t$ is zero; otherwise $\xi_t>0$ quantifies the smallest violation needed to reconcile the observed outcome.\cite{aitchison1986compositional}
\paragraph{Objective with judge-anchor prior.}
To avoid degenerate ``nearly-uniform'' vote reconstructions, we add a judge-anchor prior that pulls fan vote shares toward
judges' score shares in the absence of strong contradictory evidence. For each season $s$, we solve:
\begin{align}
\min_{\{p_{i,t}\},\{\xi_t\}} \quad
& \underbrace{\sum_{t=2}^{T_s}\sum_{i\in A_{s,t}\cap A_{s,t-1}}(p_{i,t}-p_{i,t-1})^2}_{\text{Temporal smoothness}}
+ \lambda_1 \underbrace{\sum_{t=1}^{T_s}\sum_{i\in A_{s,t}}(p_{i,t}-q^J_{i,t})^2}_{\text{Judge anchor prior}}
+ \lambda_2 \underbrace{\sum_{t=1}^{T_s}\xi_t^2}_{\text{Violation penalty}} \label{eq:taskA_obj_anchor}\\
\text{s.t.}\quad
& \sum_{i\in A_{s,t}} p_{i,t}=1,\ \ p_{i,t}\ge 0,\ \ \xi_t\ge 0, \qquad \forall t, \nonumber\\
& S^{(P)}_{e,t} \le S^{(P)}_{k,t} + \xi_t,\quad \forall t,\ \forall e\in E_{s,t},\ \forall k\in A_{s,t}\setminus E_{s,t}. \nonumber
\end{align}




\textbf{Remark (objective form).}
Equation~\eqref{eq:taskA_obj_anchor} is a convex quadratic program (QP) solved independently for each season.
The smoothness term is computed only for contestants active in consecutive weeks ($A_{s,t}\cap A_{s,t-1}$), while the anchor and slack penalties
apply to all active contestants and weeks. Hyperparameters $(\lambda_1,\lambda_2)$ are selected by replay-based diagnostics (Section~\ref{subsec:a_hyper_impl}).\cite{boyd2004convex,bubeck2015convex}

\begin{algorithm}[H]
\caption{Season-wise fan vote inference via constrained QP}
\label{alg:qp_inference}
\small  % 字号缩小，避免公式溢出
\begin{algorithmic}[1]  % [1] 表示显示行号，algpseudocode必须加
    % ---- 输入输出定义（algpseudocode标准写法）----
    \Require 
        Weekly panel for season $s$: active sets $A_{s,t}$, judges' totals $J_{i,t}$, eliminated sets $E_{s,t}$; hyperparameters $\alpha,\lambda_1,\lambda_2$.
    \Ensure 
        Point estimates $\hat p_{i,t}$ and slack diagnostics $\hat \xi_t$ for all eligible weeks.

    % ---- 算法步骤 ----
    \State Compute judges' share for each week $t$: 
    \[
    q^J_{i,t} = \frac{J_{i,t}}{\sum_{r\in A_{s,t}} J_{r,t}}
    \]
    \State Define decision variables $\{p_{i,t}\}_{i\in A_{s,t}}$ with constraints:
    \[
    p_{i,t} \ge 0, \quad \sum_{i\in A_{s,t}} p_{i,t} = 1
    \]
    \State For each week $t$, define percent-based composite score $S^{(P)}_{i,t}$ and slack variable $\xi_t$:
    \[
    S^{(P)}_{i,t} = \alpha q^{J}_{i,t} + (1-\alpha)p_{i,t}, \quad \xi_t \ge 0
    \]
    \State Add elimination-consistency constraints:
    \[
    S^{(P)}_{e,t} \le S^{(P)}_{i,t} + \xi_t \quad \forall e\in E_{s,t}, \, i\in A_{s,t}\setminus E_{s,t}
    \]
    \State Solve the quadratic program (QP) to minimize the objective function:
    \[
    \mathcal{L} = \sum_{t=2}^{T_s}\sum_{i\in A_{s,t}\cap A_{s,t-1}} (p_{i,t}-p_{i,t-1})^2 + \lambda_1\sum_{t=1}^{T_s}\sum_{i\in A_{s,t}}(p_{i,t}-q^J_{i,t})^2 + \lambda_2\sum_{t=1}^{T_s}\xi_t^2
    \]
    \State Output $\hat p_{i,t}$ (inferred vote shares) and $\hat \xi_t$ (slack diagnostics); flag weeks with large $\hat \xi_t$ as potential special-episode mismatches.
\end{algorithmic}
\end{algorithm}

\subsection{Uncertainty Quantification and Sensitivity Analysis}\label{subsec:uncertainty_qualification_and_sensitivity_analysis}
\textbf{Uncertainty is intrinsic because the inverse problem is not fully identifiable.}
We report three complementary uncertainty summaries:
(i) \textbf{Feasible intervals:} for each $(i,t)$, compute the minimum and maximum $p_{i,t}$ attainable under the same constraints (yielding an interval width as an identifiability proxy, KPI2);
(ii) \textbf{Elimination margin:} quantify how much perturbation to $\hat p_{i,t}$ is required to flip the predicted eliminated set under the replay engine, providing a week-level stability score; and
(iii) \textbf{Slack diagnostics:} $\xi_t$ indicates weeks where observed outcomes cannot be matched under the assumed mechanics without violation; these weeks are flagged for special-episode handling and for sensitivity analysis over the rule-switch season and $\alpha,\lambda_1,\lambda_2$.


\subsubsection{KPI2: Vote Identifiability (interval tightness)}\label{subsubsec:kpi2_identifiability}

\textbf{KPI2: Vote Identifiability (interval tightness).}
The KPI registry defines vote identifiability by the tightness of a feasible interval for each $(s,t,i)$ under an explicit elimination mechanism.
In Task~A, we instantiate KPI2 under the \textbf{percent-based rule} with the fan vote \textbf{share} variable $p_{i,t}$, and set $\varepsilon=10^{-6}$ to avoid singularities.
After solving the season-wise inference QP and obtaining $(\hat p,\hat\xi)$ with optimal objective value $\mathrm{OPT}_s$, we compute a \textbf{near-optimal feasible interval} for each active contestant-week pair by solving two auxiliary profile problems:
\[
p^{\min}_{i,t}=\min p_{i,t},\qquad p^{\max}_{i,t}=\max p_{i,t},
\]
subject to the original probability and elimination-consistency constraints, plus a tolerance shell that keeps solutions comparable to the fitted trajectory:
\[
\mathcal{L}(p,\xi)\le (1+\delta_{\mathrm{obj}})\mathrm{OPT}_s,\qquad \xi_\tau \le \hat\xi_\tau+\delta_\xi \ \ (\forall \tau),
\]
where $\mathcal{L}(p,\xi)$ is the Task-A objective (smoothness + judge-anchor + slack penalty).
The interval width $w_{i,t}=p^{\max}_{i,t}-p^{\min}_{i,t}$ is mapped to an identifiability score
\[
\mathrm{ID}_{i,t}=\frac{1}{w_{i,t}+\varepsilon}.
\]
We summarize $\mathrm{ID}_{i,t}$ by median/quantiles across contestants and weeks and aggregate by season and by early/mid/late phases; in implementation we compute KPI2 on selected slices (e.g., late-phase weeks or controversial cases) to control computational cost.

\paragraph{Profile optimization formulation.}
\label{para:kpi2_profile_formulation}

For a fixed season $s$, let $T_s$ be the number of weeks.
Define decision variables
\[
\mathbf{p}=\{p_{j,\tau}: j\in A_{s,\tau},\ \tau=1,\dots,T_s\},\qquad
\boldsymbol{\xi}=\{\xi_\tau:\tau=1,\dots,T_s\}.
\]
For each week $\tau$, the percent-rule combined score is
\[
S^{(P)}_{j,\tau}=\alpha q^J_{j,\tau}+(1-\alpha)p_{j,\tau}.
\]

\textbf{Base constraints.}
We impose simplex constraints for each week:\cite{aitchison1986compositional,pawlowsky2015compositional,filzmoser2018compositional}
\begin{align}
& p_{j,\tau}\ge 0,\quad \forall j\in A_{s,\tau}, \forall \tau, \label{eq:kpi2_nonneg}\\
& \sum_{j\in A_{s,\tau}}p_{j,\tau}=1,\quad \forall \tau. \label{eq:kpi2_simplex}
\end{align}
Slacks are nonnegative:
\begin{equation}
\xi_\tau\ge 0,\qquad \forall \tau. \label{eq:kpi2_slack_nonneg}
\end{equation}
Elimination-consistency constraints (same as Task~A) are
\begin{equation}
S^{(P)}_{e,\tau}\le S^{(P)}_{k,\tau}+\xi_\tau,\quad \forall \tau,\ \forall e\in E_{s,\tau},\ \forall k\in A_{s,\tau}\setminus E_{s,\tau}.
\label{eq:kpi2_elim}
\end{equation}

\textbf{Near-optimal shell (to prevent uninformative, overly wide intervals).}
Let the Task-A objective be
\begin{equation}
\mathcal{L}(\mathbf{p},\boldsymbol{\xi})
=
\sum_{\tau=2}^{T_s}\sum_{j\in A_{s,\tau}\cap A_{s,\tau-1}}(p_{j,\tau}-p_{j,\tau-1})^2
+\lambda_1\sum_{\tau=1}^{T_s}\sum_{j\in A_{s,\tau}}(p_{j,\tau}-q^J_{j,\tau})^2
+\lambda_2\sum_{\tau=1}^{T_s}\xi_\tau^2.
\label{eq:kpi2_L_def}
\end{equation}
After solving the main inference QP, we obtain $(\hat{\mathbf{p}},\hat{\boldsymbol{\xi}})$ and $\mathrm{OPT}_s=\mathcal{L}(\hat{\mathbf{p}},\hat{\boldsymbol{\xi}})$.
To compute intervals around the fitted explanation, we add:
\begin{align}
& \mathcal{L}(\mathbf{p},\boldsymbol{\xi}) \le (1+\delta_{\mathrm{obj}})\,\mathrm{OPT}_s, \label{eq:kpi2_shell_quad}\\
& \xi_\tau \le \hat{\xi}_\tau+\delta_\xi,\qquad \forall \tau. \label{eq:kpi2_slack_cap}
\end{align}
In implementation, \eqref{eq:kpi2_shell_quad} can be written as a convex SOC constraint by noting that $\mathcal{L}$ is a sum of squares:
\[
\left\|
\begin{bmatrix}
\{p_{j,\tau}-p_{j,\tau-1}\}_{\tau=2..T_s,\ j\in A_{s,\tau}\cap A_{s,\tau-1}}\\[2pt]
\{\sqrt{\lambda_1}\,(p_{j,\tau}-q^J_{j,\tau})\}_{\tau=1..T_s,\ j\in A_{s,\tau}}\\[2pt]
\{\sqrt{\lambda_2}\,\xi_\tau\}_{\tau=1..T_s}
\end{bmatrix}
\right\|_2
\le
\sqrt{(1+\delta_{\mathrm{obj}})\,\mathrm{OPT}_s}.
\]
This SOC form is directly supported by standard convex solvers (e.g., via CVX frameworks).\cite{boyd2011admm}

\textbf{Profile-min and profile-max.}
For each active pair $(i,t)$ with $i\in A_{s,t}$, define:
\begin{align}
p^{\min}_{i,t} \;=\;& \min_{\mathbf{p},\boldsymbol{\xi}}\ \ p_{i,t} \label{eq:kpi2_profile_min}\\
\text{s.t.}\ \;& \eqref{eq:kpi2_nonneg}--\eqref{eq:kpi2_elim},\ \eqref{eq:kpi2_shell_quad}--\eqref{eq:kpi2_slack_cap}, \nonumber
\end{align}
and
\begin{align}
p^{\max}_{i,t} \;=\;& \max_{\mathbf{p},\boldsymbol{\xi}}\ \ p_{i,t} \label{eq:kpi2_profile_max}\\
\text{s.t.}\ \;& \eqref{eq:kpi2_nonneg}--\eqref{eq:kpi2_elim},\ \eqref{eq:kpi2_shell_quad}--\eqref{eq:kpi2_slack_cap}. \nonumber
\end{align}
We then set the interval width $w_{i,t}=p^{\max}_{i,t}-p^{\min}_{i,t}$ and identifiability score $\mathrm{ID}_{i,t}=1/(w_{i,t}+\varepsilon)$.

\begin{algorithm}[H]
\caption{KPI2 computation via profile optimization (per season)}
\label{alg:kpi2}
\small
\begin{algorithmic}[1]  % [1] 显示行号，algpseudocode 必须加
    % ---- 输入输出（algpseudocode 标准首字母大写）----
    \Require 
        Season $s$ weekly panel: active sets $A_{s,t}$, eliminated sets $E_{s,t}$, judges' shares $q^J_{i,t}$; 
        main-model hyperparameters $\alpha,\lambda_1,\lambda_2$; 
        KPI2 params $\varepsilon,\delta_{\mathrm{obj}},\delta_\xi$; 
        a slice selector $\mathcal{S}$ (set of $(i,t)$ pairs to evaluate).
    \Ensure 
        Interval bounds $\{p^{\min}_{i,t},p^{\max}_{i,t}\}$ and identifiability scores $\{\mathrm{ID}_{i,t}\}$ for $(i,t)\in\mathcal{S}$.

    % ---- 算法核心步骤 ----
    \State \textbf{Solve main inference QP} for season $s$ to obtain $(\hat{\mathbf{p}},\hat{\boldsymbol{\xi}})$ and $\mathrm{OPT}_s=\mathcal{L}(\hat{\mathbf{p}},\hat{\boldsymbol{\xi}})$.
    \State \textbf{Define common constraints}: simplex constraints \eqref{eq:kpi2_simplex}--\eqref{eq:kpi2_elim} and slack nonnegativity \eqref{eq:kpi2_slack_nonneg}.
    \State \textbf{Add KPI2 shells}: objective shell \eqref{eq:kpi2_shell_quad} (or its SOC form) and slack cap \eqref{eq:kpi2_slack_cap}.
    \State Choose evaluation slice $\mathcal{S}$ (e.g., late-phase weeks and/or controversial cases) to control computation.
    
    % 循环结构（algpseudocode 首字母大写）
    \For{each $(i,t)\in\mathcal{S}$}
        \State Solve profile-min optimization \eqref{eq:kpi2_profile_min} to get $p^{\min}_{i,t}$.
        \State Solve profile-max optimization \eqref{eq:kpi2_profile_max} to get $p^{\max}_{i,t}$.
        \State Compute width $w_{i,t}=p^{\max}_{i,t}-p^{\min}_{i,t}$ and identifiability score $\mathrm{ID}_{i,t}=1/(w_{i,t}+\varepsilon)$.
        \State Sanity check: verify $p^{\min}_{i,t}\le \hat p_{i,t}\le p^{\max}_{i,t}$; if violated, align constraints with the main QP.
    \EndFor
    
    \State Aggregate $\mathrm{ID}_{i,t}$ by week/season/phase using medians and quantiles as defined in the KPI registry.
\end{algorithmic}
\end{algorithm}

\subsection{Hyperparameters, tuning, and implementation details}
\label{subsec:a_hyper_impl}

\textbf{Hyperparameters and defaults.}
Task~A uses $\alpha$ (judge weight in the percent rule), $\lambda_1$ (strength of the judge-anchor prior),
$\lambda_2$ (penalty on constraint violations via $\xi_t$), and the KPI2 parameters $\varepsilon,\delta_{\mathrm{obj}},\delta_\xi$.
Unless otherwise stated, we set $\alpha=0.5$ as the nominal DWTS percent weighting and treat $\alpha$ as a sensitivity knob in later stress tests.

We tune $(\lambda_1,\lambda_2)$ by replay-based diagnostics. Increasing $\lambda_2$ enforces stricter elimination-consistency
(smaller $\xi_t$) at the cost of reduced feasibility for irregular-format weeks; decreasing $\lambda_2$ allows the model to absorb
format mismatches but may underfit eliminations. Increasing $\lambda_1$ pulls inferred votes toward judges' score shares and helps
avoid near-uniform solutions, while excessively large $\lambda_1$ may suppress genuine judge--fan disagreements that are evidenced by
the observed eliminations. Practically, we tune $(\lambda_1,\lambda_2)$ on a season-by-season basis by balancing
(i) the frequency and magnitude of nonzero slacks $\hat\xi_t$,
(ii) the stability of inferred top-$k$ vote shares across adjacent weeks, and
(iii) downstream replay fidelity in Task~B.

\textbf{Special episodes and multi-elimination weeks.}
Our formulation supports (a) multi-elimination weeks by allowing $|E_{s,t}|>1$ and enforcing the same pairwise inequalities against all non-eliminated contestants, and
(b) non-elimination or irregular weeks by allowing $E_{s,t}=\varnothing$ (no elimination constraints) while still solving the smoothness objective; such weeks can also be explicitly flagged and assessed via slack and sensitivity tests.
If a season contains weeks whose outcomes cannot be reconciled with the assumed mechanics, the model will surface this via persistently large $\hat\xi_t$; these weeks are treated as format exceptions rather than forcing ad hoc re-labeling.

\textbf{Solver and numerical stability.}
The main inference problem is a convex QP and can be solved season-wise using standard QP solvers.
The KPI2 shell can be implemented either directly as a quadratic constraint or in SOC form; the latter is solver-friendly in CVX frameworks.
To avoid numerical issues, we recommend
(i) normalizing judges' scores into shares $q^J$ (already done),
(ii) using $\varepsilon=10^{-6}$ for identifiability scoring,
and (iii) keeping $\delta_{\mathrm{obj}}$ small (e.g., 1--5\%) so profile intervals remain comparable to the fitted explanation.
